\section{Conseguenze, impatti e confronto dei rischi}
\label{sec:conseguenze}

\subsection{Dal comportamento malevolo al danno}

Per valutare l'impatto reale di questi due attacchi, dobbiamo separare l'esposizione potenziale dal danno effettivo. Non basta aver scaricato il codice malevolo per considerarsi una vittima.

\textbf{Caso npm (chalk/debug).} Il codice malevolo inserito nei pacchetti intercettava le interazioni con i wallet e poteva sostituire gli indirizzi di destinazione delle transazioni~\cite{debug2025}. Perché si verificasse l'impatto, la libreria compromessa doveva finire in un bundle frontend, essere caricata nel browser ed essere eseguita durante un'operazione con un wallet. I circa due miliardi di download settimanali dei pacchetti coinvolti indicano la loro diffusione potenziale, non il numero di utenti colpiti. CSO Online riportò circa 970 dollari di fondi sottratti tracciati da Aikido~\cite{cso2025}. Oltre al danno economico, le organizzazioni hanno dovuto verificare le dipendenze e ricostruire i bundle potenzialmente esposti~\cite{wiz2025}.

\textbf{Caso AUR.} Le ricette modificate potevano avviare codice sul sistema durante l'installazione del pacchetto. L'avviso ufficiale di Arch conferma l'ondata di adozioni e aggiornamenti malevoli, ma non descrive i payload; l'analisi del pacchetto npm \texttt{atomic-lockfile} inserito nelle ricette lo classifica come \emph{infostealer}, cioè un malware che raccoglie credenziali. Tra i dati cercati figurano cookie, token di sessione, credenziali GitHub e npm e chiavi SSH~\cite{arch-aur-incident-2026,ioctl2026}. L'impatto dipende dai privilegi del processo: un'esecuzione con privilegi elevati può compromettere il sistema, mentre con privilegi utente il rischio riguarda soprattutto i dati accessibili a quell'account. Sono capacità documentate dall'analisi, non danni accertati: le fonti non indicano quante macchine siano state effettivamente compromesse.
\subsection{Confronto secondo il Cyber Risk Management}

Possiamo mappare questi due incidenti usando il framework del Cyber Risk Management. Come si vede nella Tabella~\ref{tab:rischi_comparati}, abbiamo tralasciato percentuali e punteggi di rischio arbitrari, preferendo concentrarci sui fatti e sulle dinamiche documentate, dato che non esistono metriche pubbliche omogenee per confrontare i due ecosistemi.

\begin{table}[htbp]
\centering
\footnotesize
\renewcommand{\arraystretch}{1.2}
\begin{tabular}{|p{1.4cm}|p{2cm}|p{1.6cm}|p{2cm}|p{2.3cm}|p{2.1cm}|p{2cm}|}
\hline
\textbf{Caso} & \textbf{Asset} & \textbf{Minaccia} & \textbf{Debolezza} & \textbf{Tecnica} & \textbf{Esposizione} & \textbf{Rischio / Impatto} \\ \hline
\textbf{npm} & Integrità bundle; fondi crypto degli utenti & Attore malevolo & Phishing e credenziali publisher compromesse & Versione malevola pubblicata da account legittimo & Esecuzione nel browser durante l'uso di un wallet & Frode crypto e costi di bonifica \\ \hline
\textbf{AUR} & Credenziali e dati locali & Attore malevolo & Adozione e modifica di ricette non verificate & Script e dipendenza npm malevoli & Installazione del pacchetto e avvio dei relativi hook & Furto di credenziali; possibile compromissione del sistema \\ \hline
\end{tabular}
\caption{Modellazione comparativa dei rischi per gli ecosistemi npm e AUR.}
\label{tab:rischi_comparati}
\end{table}

\textbf{Differenze e incertezze.} Quello che emerge è una forte differenza nel tipo di fiducia abusata. Nel caso npm, è saltata la fiducia nell'\emph{identità} dell'autore: l'aggiornamento era firmato dal manutentore legittimo e l'attaccante ha sfruttato l'albero delle dipendenze per colpire in modo mirato un target specifico a valle (gli utenti crypto). In AUR, invece, è stata bucata la fiducia nel \emph{contenuto} procedurale: la vittima è lo stesso sviluppatore che si fida di una ricetta orfana presa e aggiornata dalla community. 
Per quanto riguarda le incertezze, nel caso npm il danno monetario tracciato è limitato rispetto alla diffusione potenziale dei pacchetti~\cite{debug2025,cso2025}. Per AUR, invece, le centinaia di ricette segnalate descrivono modifiche al repository, ma non quantificano quante installazioni abbiano eseguito il codice né quanti sistemi siano stati compromessi~\cite{arch-aur-incident-2026,ioctl2026}. L'impatto complessivo resta quindi non quantificabile con le fonti disponibili.
